% page 496 du fac-simile (image p-495 du scan)
% bloc 160.0 x 250.6 mm ; marge gauche 8.0 mm
\pgc{-12.700mm}{0.000mm}{
\l{\cel{59}488.}
\l{}
\l{\sou{rehabilitar}. (\sou{trans.)}\cel{2}I. Ri-establisar en lua stando origi-}
\l{\cel{3}nala, en lua yuri originala (persono quan onu deklarabis}
\l{\cel{3}dekadinta de li). - II. (\sou{metaf.)}\cel{2}Ri-estimigar da lua kun-}
\l{\cel{3}homi. - DEFIRS.}
\l{}
\l{reino. Singla ek la rimeni di la brido (o di la brideto) quin}
\l{\cel{3}onu uzas por direktar kavalo. - EFIS.}
\l{}
\l{\sou{rejo}. Chefo suverena di ula stati. - II.(\sou{karto-ludo)}. Karto}
\l{\cel{3}qua reprezentas rejo en singla de la du kolori. - III.}
\l{\cel{3}(\sou{shak-ludo}). Piono precipua di la ludo, quan onu ne povas}
\l{\cel{3}prenar ; la partio esas perdita kande la rejo esas prenebla.}
\l{\cel{3}- EFIS.}
\l{}
\l{\sou{rejimo}. I. (\sou{politiko)}. Formo di guverno. II. Maniero adminis-}
\l{\cel{3}trar institucuro. - III. (\sou{medic.}) Maniero administrar sua}
\l{\cel{3}higieno individuala, sua bona-stando. - IV. (\sou{medic.}) Manie-}
\l{\cel{3}ro vivar en qua onu surveyas su, precipue pri to quon onu}
\l{\cel{3}manjas o drinkas. - VI. (\sou{tekn.)}\cel{1}Maniero ye qua eventas la}
\l{\cel{3}fluo di aquo nestagnanta. - DEFIRS.}
\l{}
\l{\sou{rekalkar}. (\sou{trans.)}\cel{2}(\sou{solduro, e}\cel{1}c.) Kunpresar, per martelo-}
\l{\cel{3}frapi, la plombo, la stano, qua salias, transpasas. - IS}
\l{}
\l{\sou{rekamar}. (\sou{trans.}) Lor texo, pasigar, sur la fondo uniforma}
\l{\cel{3}di stofo, fili qui formacas desegnuro. - IS.}
\l{}
\l{\sou{rekapitular}. (\sou{trans.}) Rezumar, kurte repetar, o memorigar,}
\l{\cel{3}per enumero preciza e kompleta di omna punti, parti o oha-}
\l{\cel{3}pitri. - DEFIS.}
\l{}
\l{\sou{reklamar}. (\sou{netrans., pri}).\cel{2}Laudar (mem exajeroze) publike}
\l{\cel{3}la komercajo, la varo, la objekto quin onu probas vendar.}
\l{\cel{3}- DFRS.}
\l{}
\l{\sou{reklamacar}.I(\sou{netrans.}) (\sou{ad, pri, por}). Protestar kontre to}
\l{\cel{3}quon onu judikas kom neyusta, por obtenar la reparo di la}
\l{\cel{3}domajo o lezo ; desagnoskar to quon onu judikas kom neyusta,}
\l{\cel{3}ne-equitatoza.II (\sou{trans}\cel{1}.) Postular (to quo debesas. III.}
\l{\cel{3}Demandar (ulo) tre insistoze. - DEF.}
\l{}
\l{\sou{rekolecar}.\cel{2}(\sou{netrans}.) I. Koncentrar su aden su, per izolar su}
\l{\cel{3}de la kozi extera. - II. Koncentrar su ad pensi di pieso.-}
\l{\cel{3}FIS.}
\l{}
\l{\sou{rekoliar}. (\sou{trans.}) Deprenar, hike ed ibe, to quo falis, disper-}
\l{\cel{3}site, sur sulo e retro-pozar lu a lua plaso. - F.}
\l{}
\l{\sou{rekoltar}. (\sou{trans.)}\cel{2}Rekoliar la produkturi da la tero kultivi-}
\l{\cel{3}ta. - (\sou{anke metaf}.)FI.}
\l{}
\l{rekomendar. (trans.)\cel{2}I. (\sou{pri persono)}.Indikar (lu) partikula-}
\l{\cel{3}re kom esonta la objekto di la bonvolo, di la protekto, da}
\l{\cel{3}ulu. II. (pri letro, pako, e c.) (En posto-kontoro), deman-}
\l{\cel{3}\sou{dar}\cel{1}ke la letro-livristo donez lu a la destinario ipsa, ne-}
\l{\cel{3}mediate - po pago di\cel{1}\sou{afr}anko-suplemento. - DEFIRS}
}
